% Compiled with: latexmk -pdf -interaction=nonstopmode main.tex
\documentclass[12pt]{article}
\usepackage[margin=1in]{geometry}
\usepackage{amsmath,amssymb,booktabs,graphicx,setspace,caption}
\usepackage[longnamesfirst]{natbib}
\usepackage[hidelinks]{hyperref}
\onehalfspacing
\bibliographystyle{jf} % swap for target journal; jf.bst not in texlive-core — fall back:
% \bibliographystyle{plainnat}

\title{Working Title\thanks{Author affiliation and acknowledgments. Draft — do not circulate.}}
\author{}
\date{\today}

\begin{document}
\maketitle

\begin{abstract}
\noindent 100--150 words: question, setting, identification in one clause, headline magnitude with units, mechanism claim no stronger than the evidence.
\end{abstract}

\section{Introduction}
% Contribution stated by page 2. Lead with the raw-data fact, then the baseline estimate.

\section{Related literature}
% Built from docs/LIT_NOTES.md — COMPETES papers addressed head-on.

\section{Data}

\section{Empirical strategy}
% The estimating equation, the identifying assumption in one sentence, and the threat/defense pairs.

\section{Results}
% Figures from paper/figures/, tables generated into paper/tables/ — never hand-typed.
% \begin{figure}[!htbp]\centering\includegraphics[width=.8\textwidth]{figures/raw_fact.pdf}\caption{}\label{fig:raw}\end{figure}
% \input{tables/baseline.tex}

\section{Robustness}
% Includes nulls, failed placebos, and open flags per STANDING_RULES.md rule 15.

\section{Conclusion}

\bibliography{references}

\appendix
\section{Additional results}

\end{document}
